\sectionheading{5.10 চতুর্থ প্রজন্মের ভাষা (Fourth Generation Language — 4GL)}

কম্পিউটারের ব্যবহার আরও সহজ করার উদ্দেশ্যে এমন ভাষার বিকাশ ঘটে,
যেখানে Programmer-কে প্রতিটি কাজের বিস্তারিত ধাপ লিখতে হয় না;
বরং তুলনামূলক সহজ ও স্বাভাবিক ভাষার কাছাকাছি নির্দেশ দিয়ে
প্রয়োজনীয় কাজ সম্পন্ন করা যায়।

এ ধরনের ভাষাকে
\textbf{Fourth Generation Language (4GL)}
বা \textbf{চতুর্থ প্রজন্মের ভাষা} বলা হয়।

4GL-কে অনেক ক্ষেত্রে \textbf{অতি উচ্চস্তরের ভাষা}
হিসেবেও উল্লেখ করা হয়।

%================================================
% Definition
%================================================
\booksubheading{চতুর্থ প্রজন্মের ভাষা কাকে বলে?}

\begin{definitionbox}

\textbf{\color{darkTeal}সংজ্ঞা}

যে Programming Language-এ তুলনামূলক কম ও সহজ নির্দেশ ব্যবহার করে
Database Query, Report তৈরি, Data পরিচালনা এবং বিভিন্ন Application-এর
কাজ সম্পন্ন করা যায়, তাকে
\textbf{চতুর্থ প্রজন্মের ভাষা
(Fourth Generation Language — 4GL)} বলে।

\end{definitionbox}

\textbf{সহজভাবে মনে রাখবে—}

\begin{center}

\textbf{
কম নির্দেশ
\(\;+\;\)
সহজ ব্যবহার
\(\;+\;\)
দ্রুত কাজ
\(\;=\;\)
4GL
}

\end{center}

%================================================
% Main Concept
%================================================
\booksubheading{4GL-এর মূল ধারণা}

আগের প্রজন্মের অনেক Programming Language-এ Programmer-কে
একটি কাজ \textbf{কীভাবে (How)} সম্পন্ন হবে তার ধাপ
তুলনামূলক বিস্তারিতভাবে লিখতে হয়।

4GL-এ অনেক ক্ষেত্রে Programmer মূলত
\textbf{কী ফলাফল (What)} প্রয়োজন তা জানায় এবং প্রয়োজনীয়
অভ্যন্তরীণ কাজ System নিজে সম্পন্ন করে।

\begin{center}

\fcolorbox{darkBlue}{softBlue}{
    \parbox[c][10mm][c]{0.34\linewidth}{
        \centering
        \textbf{3GL}\\
        How — কীভাবে কাজ হবে
    }
}
\hspace{7mm}
\fcolorbox{mainTeal}{softGreen}{
    \parbox[c][10mm][c]{0.34\linewidth}{
        \centering
        \textbf{4GL}\\
        What — কী ফল প্রয়োজন
    }
}

\end{center}

\begin{examplebox}{সহজ উদাহরণ}

Database থেকে নির্দিষ্ট Data বের করার জন্য SQL-এ একটি Query
লেখা যায়।

Programmer-কে Database-এর প্রতিটি Record কোন ধাপে বা কীভাবে
Search করা হবে তা আলাদাভাবে লিখতে হয় না।

অর্থাৎ Programmer মূলত \textbf{কোন Data প্রয়োজন} তা জানায়।

\end{examplebox}

%================================================
% Application Generator
%================================================
\booksubheading{4GL-কে Application Generator বলা হয় কেন?}

4GL ব্যবহার করে অনেক ক্ষেত্রে তুলনামূলক কম Coding-এর মাধ্যমে
Form, Query, Report এবং Data Processing-এর কাজ দ্রুত তৈরি করা যায়।

এ কারণে কিছু 4GL Tool-কে
\textbf{Application Generator} হিসেবেও উল্লেখ করা হয়।

\begin{center}

\textbf{
কম Coding
\(\longrightarrow\)
কম সময় ও শ্রম
\(\longrightarrow\)
দ্রুত Application Development
}

\end{center}

%================================================
% Features
%================================================
\booksubheading{4GL-এর প্রধান বৈশিষ্ট্য}

\begin{tabularx}{\linewidth}{|>{\bfseries}p{0.28\linewidth}|Y|}
\hline

\rowcolor{softBlue}
\multicolumn{1}{|c|}{\textbf{বৈশিষ্ট্য}}
&
\multicolumn{1}{c|}{\textbf{সহজ ব্যাখ্যা}}
\\
\hline

অতি উচ্চস্তরের ভাষা
&
মানুষের ব্যবহার ও Problem Solving-এর তুলনামূলক কাছাকাছি।
\\
\hline

Natural Language-এর কাছাকাছি
&
নির্দেশ তুলনামূলক সহজ ও বোধগম্য।
\\
\hline

কম Coding
&
একই কাজের জন্য তুলনামূলক কম নির্দেশ প্রয়োজন হয়।
\\
\hline

Query সুবিধা
&
Database থেকে প্রয়োজনীয় Data অনুসন্ধান ও নির্বাচন করা যায়।
\\
\hline

Report Generation
&
Data থেকে সহজে Report তৈরি করা যায়।
\\
\hline

দ্রুত Development
&
Application বা Software তৈরিতে সময় কম লাগে।
\\
\hline

কম শ্রম
&
Programming-এর প্রয়োজনীয় শ্রম তুলনামূলক কম হয়।
\\
\hline

User-oriented
&
ব্যবহারকারীর প্রয়োজনীয় Result দ্রুত পাওয়ার ওপর বেশি গুরুত্ব দেওয়া হয়।
\\
\hline

\end{tabularx}

%================================================
% Uses
%================================================
\booksubheading{4GL-এর ব্যবহার}

4GL বিশেষভাবে Data ও Business-oriented কাজের ক্ষেত্রে
ব্যবহার করা হয়। যেমন—

\begin{compacttwocolpoints}

\item Database Query তৈরিতে।

\item Data সংরক্ষণ ও অনুসন্ধানে।

\item Report তৈরিতে।

\item Data পরিচালনা ও আদান-প্রদানে।

\item Business Application তৈরিতে।

\item দ্রুত Application Development-এ।

\end{compacttwocolpoints}

%================================================
% Query Language
%================================================
\booksubheading{Query Language কী?}

\begin{definitionbox}

\textbf{\color{darkTeal}Query Language}

Database থেকে প্রয়োজনীয় Data অনুসন্ধান, নির্বাচন বা পরিবর্তনের
জন্য ব্যবহৃত ভাষাকে সাধারণভাবে
\textbf{Query Language} বলা হয়।

\end{definitionbox}

4GL-এর আলোচনায়
\textbf{SQL (Structured Query Language)}
একটি গুরুত্বপূর্ণ উদাহরণ।

%================================================
% SQL Example
%================================================
\begin{examplebox}{SQL Query-এর সহজ উদাহরণ}

ধরা যাক, \textenglish{\ttfamily student} নামের Table-এর
সব Data দেখতে হবে।

\begin{codebox}

\begin{center}
{\textenglish{\ttfamily SELECT * FROM student;}}
\end{center}

\end{codebox}

এখানে Programmer মূলত কোন Data প্রয়োজন তা Query-এর মাধ্যমে
Database System-কে জানাচ্ছে।

\end{examplebox}

%================================================
% Examples
%================================================
\booksubheading{4GL-এর উদাহরণ}

4GL-এর সবচেয়ে পরিচিত উদাহরণ হিসেবে—

\begin{center}

\textbf{SQL (Structured Query Language)}

\end{center}

উল্লেখ করা যায়।

কিছু পাঠ্যবইয়ে 4GL-এর উদাহরণের আলোচনায় Oracle-এর নামও পাওয়া যায়।
তবে \textbf{Oracle নিজে Java বা Python-এর মতো General Purpose
Programming Language নয়}; Oracle হলো একটি
\textbf{RDBMS/Database Platform}।

Oracle Database Environment-এ SQL ও PL/SQL ব্যবহার করা হয়।

\begin{rememberbox}

\textbf{\color{deepOrange}সতর্কতা}

পরীক্ষায়—

\textbf{SQL = 4GL-এর গুরুত্বপূর্ণ উদাহরণ}

লেখা যায়।

কিন্তু—

\textbf{Oracle = General Purpose Programming Language}

লেখা ঠিক নয়।

Oracle মূলত একটি \textbf{RDBMS / Database Platform}।

\end{rememberbox}

%================================================
% Advantages
%================================================
\booksubheading{4GL-এর সুবিধা}

\begin{compacttwocolpoints}

\item Programming তুলনামূলক সহজ।

\item কম Code লিখে বেশি কাজ করা যায়।

\item Application দ্রুত তৈরি করা যায়।

\item Database Query তৈরি সহজ।

\item Report Generation সহজ।

\item Software Development-এ সময় কম লাগে।

\item Programming-এর শ্রম তুলনামূলক কমে।

\item Programmer-এর Productivity বাড়ে।

\item Business ও Database Application তৈরিতে সুবিধাজনক।

\end{compacttwocolpoints}

%================================================
% Limitations
%================================================
\booksubheading{4GL-এর কিছু সীমাবদ্ধতা}

\begin{itemize}

\item সব ধরনের System-level Programming-এর জন্য উপযোগী নয়।

\item Hardware-এর ওপর Low Level Language-এর মতো সরাসরি
নিয়ন্ত্রণ পাওয়া যায় না।

\item নির্দিষ্ট 4GL Tool বা Database Environment-এর ওপর
নির্ভরতা থাকতে পারে।

\item অত্যন্ত বিশেষ বা Hardware-oriented কাজের জন্য
General Purpose বা Low Level Language প্রয়োজন হতে পারে।

\end{itemize}

%================================================
% 3GL vs 4GL
%================================================
\booksubheading{3GL ও 4GL-এর পার্থক্য}

\begin{tabularx}{\linewidth}{|>{\bfseries}p{0.21\linewidth}|Y|Y|}
\hline

\rowcolor{softBlue}
\multicolumn{1}{|c|}{\textbf{বিষয়}}
&
\multicolumn{1}{c|}{\textbf{3GL}}
&
\multicolumn{1}{c|}{\textbf{4GL}}
\\
\hline

মূল গুরুত্ব
&
কাজ কীভাবে সম্পন্ন হবে
&
কী Result প্রয়োজন
\\
\hline

Instruction
&
তুলনামূলক বেশি
&
তুলনামূলক কম
\\
\hline

Code-এর পরিমাণ
&
বেশি হতে পারে
&
তুলনামূলক কম
\\
\hline

Development
&
তুলনামূলক বেশি সময় লাগতে পারে
&
দ্রুত Development সম্ভব
\\
\hline

ব্যবহার
&
General Programming
&
Query, Report ও Data/Business Application
\\
\hline

Programming জ্ঞান
&
তুলনামূলক বেশি প্রয়োজন
&
অনেক কাজ সহজভাবে করা যায়
\\
\hline

উদাহরণ
&
C, FORTRAN, BASIC
&
SQL
\\
\hline

\end{tabularx}

%================================================
% Generation
%================================================
\booksubheading{4GL ও ভাষার প্রজন্ম}

Fourth Generation Language হলো Programming Language-এর
\textbf{চতুর্থ প্রজন্ম}।

\begin{center}

\textbf{
4GL
\(\;=\;\)
Fourth Generation Language
}

\end{center}

4GL-এর বিকাশ 1970-এর দশক থেকে উল্লেখযোগ্যভাবে শুরু হয় বলে
প্রচলিত পাঠ্যবইয়ে আলোচনা করা হয়।

এর উদ্দেশ্য ছিল Programming আরও সহজ করা এবং কম Coding-এর মাধ্যমে
দ্রুত প্রয়োজনীয় Result বা Application তৈরি করা।

%================================================
% Final Reminder
%================================================
\begin{rememberbox}

\textbf{\color{deepOrange}মনে রাখবে}

\begin{itemize}

\item 4GL-এর পূর্ণরূপ
\textbf{Fourth Generation Language}।

\item 4GL-কে \textbf{অতি উচ্চস্তরের ভাষা} হিসেবেও উল্লেখ করা হয়।

\item 4GL-এর বিকাশ 1970-এর দশক থেকে আলোচিত।

\item 4GL-এ তুলনামূলক
\textbf{কম Code-এ বেশি কাজ} করা যায়।

\item 4GL-এ অনেক ক্ষেত্রে
\textbf{How-এর চেয়ে What}-এর ওপর বেশি গুরুত্ব দেওয়া হয়।

\item 4GL বিশেষভাবে
\textbf{Database Query, Report ও Business Application}-এ উপযোগী।

\item কিছু 4GL Tool-কে
\textbf{Application Generator} বলা হয়।

\item \textbf{SQL} 4GL-এর একটি গুরুত্বপূর্ণ উদাহরণ।

\item SQL-এর পূর্ণরূপ
\textbf{Structured Query Language}।

\item 4GL ব্যবহারে Software Development-এর
\textbf{সময় ও শ্রম কমতে পারে}।

\end{itemize}

\end{rememberbox}

\small

%================================================
% 5.11 Translator Program
%================================================
